\documentclass[12pt,a4paper]{article}
\usepackage[utf8]{inputenc}
\usepackage[spanish]{babel}
\usepackage{amsmath}
\usepackage{amsfonts}
\usepackage{amssymb}
\usepackage{tikz}
\usetikzlibrary{arrows.meta, positioning, calc}
\usepackage[left=2cm,right=2cm,top=2cm,bottom=2cm]{geometry}

\begin{document}

\begin{center}
    \Large \textbf{Sistemas Distribuidos} \\
    \large \textbf{Ejercicios Sincronización Relojes}
\end{center}

\vspace{0.5cm}

\textbf{1- Un cliente sincroniza su reloj usando el algoritmo de Cristian. Si se tienen los siguientes datos:}
\begin{itemize}
    \item $t_0 = 10:00:00$ (envío)
    \item $t_1 = 10:00:06$ (recepción)
    \item tiempo del servidor: $C_s = T_s = 10:00:03$
\end{itemize}
\textbf{¿A qué hora debe ajustar su reloj el cliente?}

\vspace{0.3cm}
\textbf{Solución:}
Aplicamos la fórmula del algoritmo de Cristian: 
$$ C = C_s + \frac{t_1 - t_0}{2} $$

Calculamos el tiempo de viaje redondo (RTT):
$$ RTT = t_1 - t_0 = 10:00:06 - 10:00:00 = 6 \text{ segundos} $$

Calculamos el tiempo de ajuste (latencia de un solo sentido):
$$ \frac{RTT}{2} = \frac{6}{2} = 3 \text{ segundos} $$

Ajustamos el reloj del cliente:
$$ C = 10:00:03 + 00:00:03 = 10:00:06 $$
\textbf{El cliente debe ajustar su reloj a las 10:00:06.}

\vspace{0.8cm}

\textbf{2- Si en el algoritmo de Cristian existe latencia asimétrica, es decir $t_0=12:00:00$, $t_1=12:00:10$ y $C_s=T_s=12:00:04$. Pero, la duración de ida es de 2 s y la de vuelta de 8 s. }\\
\textbf{¿Qué valor del tiempo calcula el algoritmo de Cristian? ¿Cuál sería el tiempo real correcto?}

\vspace{0.3cm}
\textbf{Solución:}
\\
\textit{a) Valor que calcula el algoritmo de Cristian:} \\
El algoritmo de Cristian asume latencia simétrica. 
$$ RTT = 12:00:10 - 12:00:00 = 10 \text{ segundos} $$
$$ \text{Ajuste} = \frac{10}{2} = 5 \text{ segundos} $$
$$ C = 12:00:04 + 00:00:05 = \textbf{12:00:09} $$

\textit{b) Tiempo real correcto:} \\
Como sabemos que el viaje de vuelta tarda exactamente 8 segundos, el tiempo real en el cliente cuando recibe el mensaje debería ser el tiempo del servidor más esos 8 segundos de viaje:
$$ \text{Tiempo real} = C_s + \text{vuelta} = 12:00:04 + 00:00:08 = \textbf{12:00:12} $$

\vspace{0.8cm}

\textbf{3- Usando el algoritmo de sincronización de Berkeley dado en clase y los siguientes datos:}
\begin{center}
\begin{tabular}{|c|c|}
\hline
\textbf{Nodo} & \textbf{Hora} \\
\hline
Servidor & 10:00:00 \\
\hline
A & 10:00:04 \\
\hline
B & 09:59:58 \\
\hline
C & 10:00:02 \\
\hline
\end{tabular}
\end{center}

\vspace{0.3cm}
\textbf{Solución paso a paso:}

\textbf{Pasos 1, 2 y 3:} Selección del servidor. El servidor envía solicitudes y recibe la hora actual de los clientes.
\begin{center}
\shorthandoff{>}
\begin{tikzpicture}[x=2cm, y=2cm, 
    node box/.style={rectangle, draw=black, thick, minimum width=1.5cm, minimum height=1.5cm, fill=gray!10, text centered}]
    
    % Servidor
    \node[node box] (S) at (0, 0) {\Large S};
    \node[above=0.1cm of S, text=blue, font=\bfseries] {10:00:00};
    
    % Clientes
    \node[node box] (A) at (-2, -1.5) {\Large A};
    \node[above=0.1cm of A, text=blue, font=\bfseries] {10:00:04};
    
    \node[node box] (B) at (0, -1.5) {\Large B};
    \node[above=0.1cm of B, text=blue, font=\bfseries] {09:59:58};
    
    \node[node box] (C) at (2, -1.5) {\Large C};
    \node[above=0.1cm of C, text=blue, font=\bfseries] {10:00:02};

    % Flechas de los clientes al servidor (enviando su tiempo)
    \draw[->, thick] (A.north) to[bend left=15] (S.west);
    \draw[->, thick] (B.north) -- (S.south);
    \draw[->, thick] (C.north) to[bend right=15] (S.east);
\end{tikzpicture}
\shorthandon{>}
\end{center}

\textbf{Paso 4 y parte 1 del ejercicio (Cálculo de diferencias y Promedio):} \\
El servidor calcula la diferencia de tiempo entre los relojes de los clientes y el suyo ($T_{cliente} - T_{servidor}$):
\begin{center}
\shorthandoff{>}
\begin{tikzpicture}[x=2.5cm, y=2cm, 
    node box/.style={rectangle, draw=black, thick, minimum width=1.5cm, minimum height=1.5cm, fill=gray!10, text centered}]
    
    \node[node box] (S) at (0, 0) {\Large S};
    \node[node box] (A) at (-2, -1.5) {\Large A};
    \node[node box] (B) at (0, -1.5) {\Large B};
    \node[node box] (C) at (2, -1.5) {\Large C};

    % Flechas S a clientes con las diferencias
    \draw[->, thick] (S.west) to[bend right=15] node[midway, left, text=blue, align=center] {10:00:04 - 10:00:00 \\ \textbf{+4 s}} (A.north);
    
    \draw[->, thick] (S.south) -- node[midway, right, text=blue, align=center] {09:59:58 - 10:00:00 \\ \textbf{-2 s}} (B.north);
    
    \draw[->, thick] (S.east) to[bend left=15] node[midway, right, text=blue, align=center] {10:00:02 - 10:00:00 \\ \textbf{+2 s}} (C.north);
\end{tikzpicture}
\shorthandon{>}
\end{center}

Una vez obtenidas las diferencias, ahora sí procedemos a calcular el \textbf{promedio} (incluyendo la del propio servidor que es 0):
$$ \text{Promedio} = \frac{+4 - 2 + 2 + 0}{4} = \frac{4}{4} = \textbf{+1 segundo} $$
Esto significa que la nueva hora consensuada de la red será la hora actual del servidor más 1 segundo: \textbf{10:00:01}.

\vspace{0.3cm}
\textbf{Paso 5 y parte 2 del ejercicio (Determinar el ajuste para cada nodo):} \\
El servidor calcula cuánto debe ajustarse cada nodo para llegar a la hora promedio (10:00:01) y les envía el ajuste pertinente ($\text{Hora Promedio} - \text{Hora Nodo}$):
\begin{itemize}
    \item \textbf{Servidor S:} ajusta $10:00:01 - 10:00:00 = \textbf{+1 s}$
    \item \textbf{Nodo A:} ajusta $10:00:01 - 10:00:04 = \textbf{-3 s}$
    \item \textbf{Nodo B:} ajusta $10:00:01 - 09:59:58 = \textbf{+3 s}$
    \item \textbf{Nodo C:} ajusta $10:00:01 - 10:00:02 = \textbf{-1 s}$
\end{itemize}

\begin{center}
\shorthandoff{>}
\begin{tikzpicture}[x=2cm, y=2cm, 
    node box/.style={rectangle, draw=black, thick, minimum width=1.5cm, minimum height=1.5cm, fill=gray!10, text centered}]
    
    \node[node box] (S) at (0, 0) {\Large S};
    \node[node box] (A) at (-2, -1.5) {\Large A};
    \node[node box] (B) at (0, -1.5) {\Large B};
    \node[node box] (C) at (2, -1.5) {\Large C};

    % Flechas de ajustes
    \draw[->, thick] (S.west) to[bend right=15] node[midway, left, text=blue] {\textbf{-3 s}} (A.north);
    \node[below=0.1cm of A, text=red, font=\bfseries] {10:00:01};

    \draw[->, thick] (S.south) -- node[midway, right, text=blue] {\textbf{+3 s}} (B.north);
    \node[below=0.1cm of B, text=red, font=\bfseries] {10:00:01};

    \draw[->, thick] (S.east) to[bend left=15] node[midway, right, text=blue] {\textbf{-1 s}} (C.north);
    \node[below=0.1cm of C, text=red, font=\bfseries] {10:00:01};
    
    \node[above=0.1cm of S, text=red, font=\bfseries] {10:00:01};
\end{tikzpicture}
\shorthandon{>}
\end{center}
\newpage

\textbf{4- Dibujar un escenario con tres procesos como lo visto en clase. Considerar que los procesos tienen los siguientes eventos (A,B,C,D,E,F):} \\
P1: $A \rightarrow B \rightarrow C$ \\
P2: $D \rightarrow E$ \\
P3: $F$ \\
\textbf{Considerar que A envía a D, D envía a B, E envía a F. Además, asignar los tiempos o relojes utilizando el algoritmo de Lamport y muestre una secuencia posible de eventos.}

\vspace{0.3cm}
\textbf{Solución:} \\
Aplicando el algoritmo de Lamport ($Cp = \max(Cp, Cm) + 1$):

\begin{center}
\shorthandoff{>}
\begin{tikzpicture}[x=1.5cm, y=1.5cm]
    % Líneas de proceso
    \draw[thick, ->] (0, 3) -- (8, 3) node[right] {P1};
    \draw[thick, ->] (0, 1.5) -- (8, 1.5) node[right] {P2};
    \draw[thick, ->] (0, 0) -- (8, 0) node[right] {P3};
    
    % Eventos
    \node[circle, draw=red, thick, inner sep=2pt, fill=white] (A) at (1, 3) {A};
    \node[circle, draw=red, thick, inner sep=2pt, fill=white] (D) at (2.5, 1.5) {D};
    \node[circle, draw=red, thick, inner sep=2pt, fill=white] (B) at (4, 3) {B};
    \node[circle, draw=red, thick, inner sep=2pt, fill=white] (E) at (5, 1.5) {E};
    \node[circle, draw=red, thick, inner sep=2pt, fill=white] (C) at (6, 3) {C};
    \node[circle, draw=red, thick, inner sep=2pt, fill=white] (F) at (6.5, 0) {F};

    % Mensajes
    \draw[->, thick] (A) -- (D);
    \draw[->, thick] (D) -- (B);
    \draw[->, thick] (E) -- (F);

    % Relojes lógicos Lamport (cajas azules)
    \node[draw=blue, thick, above=0.15cm of A] {1};
    \node[draw=blue, thick, below=0.15cm of D] {2};
    \node[draw=blue, thick, above=0.15cm of B] {3};
    \node[draw=blue, thick, below=0.15cm of E] {3};
    \node[draw=blue, thick, above=0.15cm of C] {4};
    \node[draw=blue, thick, below=0.15cm of F] {4};
    
    % Ticks en P1
    \foreach \x in {0.5, 1.5, ..., 7.5} \draw (\x, 3.1) -- (\x, 2.9);
    % Ticks en P2
    \foreach \x in {0.5, 1.5, ..., 7.5} \draw (\x, 1.6) -- (\x, 1.4);
    % Ticks en P3
    \foreach \x in {0.5, 1.5, ..., 7.5} \draw (\x, 0.1) -- (\x, -0.1);
\end{tikzpicture}
\shorthandon{>}
\end{center}

\textbf{Secuencia posible de eventos:} A $\rightarrow$ D $\rightarrow$ B $\rightarrow$ E $\rightarrow$ C $\rightarrow$ F (ordenados por reloj lógico ascendente, resolviendo empates arbitrariamente o por ID de proceso).

\vspace{0.8cm}

\textbf{5- Resolver el problema 4 con relojes vectoriales. Además, indique si hay procesos concurrentes y justificar la respuesta.}

\vspace{0.3cm}
\textbf{Solución:} \\
Aplicando el algoritmo de Relojes Vectoriales ($Vp[i] = \max(Vp[i], Vm[i])$ e incrementando la componente propia en cada evento):

\begin{center}
\shorthandoff{>}
\begin{tikzpicture}[x=1.8cm, y=1.8cm]
    % Líneas de proceso
    \draw[thick, ->] (0, 3) -- (7, 3) node[right] {P1};
    \draw[thick, ->] (0, 1.5) -- (7, 1.5) node[right] {P2};
    \draw[thick, ->] (0, 0) -- (7, 0) node[right] {P3};
    
    % Eventos
    \node[circle, draw=red, thick, inner sep=2pt, fill=white] (A) at (1, 3) {A};
    \node[circle, draw=red, thick, inner sep=2pt, fill=white] (D) at (2.5, 1.5) {D};
    \node[circle, draw=red, thick, inner sep=2pt, fill=white] (B) at (4, 3) {B};
    \node[circle, draw=red, thick, inner sep=2pt, fill=white] (E) at (5, 1.5) {E};
    \node[circle, draw=red, thick, inner sep=2pt, fill=white] (C) at (5.5, 3) {C};
    \node[circle, draw=red, thick, inner sep=2pt, fill=white] (F) at (6, 0) {F};

    % Mensajes
    \draw[->, thick] (A) -- (D);
    \draw[->, thick] (D) -- (B);
    \draw[->, thick] (E) -- (F);

    % Relojes Vectoriales
    \node[above=0.1cm of A, text=blue, font=\small] {[1,0,0]};
    \node[below=0.1cm of D, text=blue, font=\small] {[1,1,0]};
    \node[above=0.1cm of B, text=blue, font=\small] {[2,1,0]};
    \node[below=0.1cm of E, text=blue, font=\small] {[1,2,0]};
    \node[above=0.1cm of C, text=blue, font=\small] {[3,1,0]};
    \node[below=0.1cm of F, text=blue, font=\small] {[1,2,1]};
    
    % Ticks
    \foreach \x in {0.5, 1.5, ..., 6.5} \draw (\x, 3.1) -- (\x, 2.9);
    \foreach \x in {0.5, 1.5, ..., 6.5} \draw (\x, 1.6) -- (\x, 1.4);
    \foreach \x in {0.5, 1.5, ..., 6.5} \draw (\x, 0.1) -- (\x, -0.1);
\end{tikzpicture}
\shorthandon{>}
\end{center}

\textbf{Procesos concurrentes ($x \parallel y$):} \\
Dos eventos son concurrentes si no se cumple que $V(x) < V(y)$ ni que $V(y) < V(x)$. En este caso, encontramos las siguientes concurrencias:
\begin{itemize}
    \item \textbf{C y E ($C \parallel E$):} Comparamos $V(C) = [3,1,0]$ y $V(E) = [1,2,0]$. Como la primera componente de C es mayor ($3 > 1$) pero la segunda componente de E es mayor ($2 > 1$), ninguno precede al otro.
    \item \textbf{C y F ($C \parallel F$):} Comparamos $V(C) = [3,1,0]$ y $V(F) = [1,2,1]$. Como la primera componente de C es mayor ($3 > 1$) pero la segunda y tercera de F son mayores ($2 > 1$ y $1 > 0$), son concurrentes.
    \item \textbf{B y E ($B \parallel E$):} Comparamos $V(B) = [2,1,0]$ y $V(E) = [1,2,0]$. La primera de B es mayor ($2 > 1$) pero la segunda de E es mayor ($2 > 1$), por lo tanto son concurrentes.
\end{itemize}

\end{document}